%% kapa: exam=1 round=24 year=2013 period=2 form=A subject=accounting subject_ko=회계학 seq=10 number=10
%% source: Q-Net 감정평가사 기출문제 (cst00309) artlSeq=5073802 file=extracted/1008/5073802/2013년 제24회 감정평가사 1차시험 기출문제/2교시(A형).hwp
%% conversion: kordoc 4.24.51 -> tools/gen1.py (auto); review=proofread-sonnet; flags=none
\begin{question}{10}
무형자산에 관한 설명으로 옳지 않은 것은?
\begin{choices}
  \item 무형자산을 최초로 인식할 때에는 공정가치로 측정한다.
  \item 최초에 비용으로 인식한 무형항목에 대한 지출은 그 이후에 무형자산의 원가로 인식할 수 없다.
  \item 자산에서 발생하는 미래경제적효익이 기업에 유입될 가능성이 높고 자산의 원가를 신뢰성 있게 측정할 수 있을 때에만 무형자산을 인식한다.
  \item 자산을 사용가능한 상태로 만드는데 직접적으로 발생하는 종업원 급여와 같은 직접 관련되는 원가는 무형자산의 원가에 포함한다.
  \item 새로운 지역에서 또는 새로운 계층의 고객을 대상으로 사업을 수행하는데서 발생하는 원가 등은 무형자산 원가에 포함하지 않는다.
\end{choices}
\end{question}
\answer{1}
